% ---------
% --------
%!TEX TS-program = pdflatex
%!TEX encoding = UTF-8 Unicode
%!TEX spellcheck = English (Aspell)

\documentclass[11pt]{article}
\usepackage{jeffe,handout,graphicx,varwidth}
\usepackage{fourier-orns}
\usepackage{mathtools}
\usepackage{braket}


\def\Sym#1{\texttt{\color{BrickRed}#1}}
\allowdisplaybreaks

\begin{document}

\begin{center}
\Large\textbf{CSCI 432/532, Spring 2026}%
\\
\LARGE\textbf{Homework 11}%
\\[0.5ex]
\large Due Monday, April 27 at 9:00am
\end{center}

\bigskip
\hrule
\bigskip

\subsection*{Submission Requirements}
\begin{itemize}
    \item Type or clearly hand-write your solutions into a PDF format so that they are legible and professional. Submit your PDF on Canvas or Gradescope. 
	\item Put each \textbf{sub}problem on a separate page and select the corresponding problem via the Gradescope interface.
    \item Do not submit your first draft. Type or clearly re-write your solutions for your final submission. If your submission is not legible, we will ask you to resubmit.
\end{itemize}

\subsection*{Collaboration and Resources}

The \emph{best} resources for completing the homework are (in no particular order):
\begin{itemize}
	\item Your own notes from lectures and problem sessions, and/or lecture recordings.
	\item The lecture notes linked from the course website for the lecture day.
	\item The questions channel on Discord.
	\item Office hours.
	\item Discussions with classmates.
\end{itemize}
You may also access almost any resource in preparing your solution to the
homework. The exception is that you may not seek answers to the problems on the
homework directly, such as by pasting them into a GenAI tool, searching for
solutions on a search engine, looking for them on Chegg or similar websites, or
asking another person for the solution.

Additionally, you \EMPH{must}
\begin{itemize}
    \item Write your solutions in your own words---this means that you should never be \emph{copying}
    anything of substance from any source, and
    \item credit every resource you use, including GenAI tools, by providing links or full chat transcripts.
    If you use the provided LaTeX template, you can use the \texttt{Sources}
    environment for this. Ask if you need help!
\end{itemize}


\begin{boxquote}{0.9}
Remember, if you choose to use GenAI tools, not only must you provide a full
transcript of your conversation (or a link to one), you must also use the
following prompt (or some variation of it), and provide a full transcript of
the conversation.
\end{boxquote}  

Prompt:
\begin{boxquote}{0.9}
You are acting as a teaching assistant for an advanced undergraduate/graduate
algorithms and computer science theory course.

Your role is not to provide full solutions or final answers.

Instead, you should act like we are having a conversation. Ask me a single
question and let me respond. Avoid asking me many questions at once or giving
me a lot of information at one time. Guide me using hints, leading questions,
partial insights, and sanity checks.  Help me debug my own proofs, algorithms,
or reductions. Do not give a complete solution or full proofs. Be patient with
me. Don't give me details so that I can move on to the next part of a problem.
Make sure that I understood before moving on.

The course uses Jeff Erickson's Models of Computation lecture notes, so you should
guide me to answer problems using the definitions, notation, and style from those notes.

Assume I am a serious student trying to learn, not to shortcut the assignment.
I am attaching the assignment, so you should refuse to provide a complete solution
to the problems listed. Of course, you can walk me through the "Solved
Problems" listed at the end if I ask.
\end{boxquote}  

\newpage
%----------------------------------------------------------------------

\headers{CSCI 432/532}{Homework 11 (due April 27)}{Spring 2026}

%----------------------------------------------------------------------
\begin{enumerate}


\item Do you think that all classical \textbf{reversible} operations also valid quantum operations? If not,
give a countereample. If yes, explain why. Any reasonable attempt will get credit, so try
to think about it yourself rather than looking up an answer!

\item Prove that the following quantum algorithm on two qubits is unitary by showing that
it preserves the unitary property on $\ket{00}$, $\ket{01}$, $\ket{10}$, and $\ket{11}$.

\begin{algo}
	\textul{$\mathsc{DoubleHadamard}$:}\+\\
	initialize two qubits $A$, $B$\\
    put $A$ and $B$ into any basic state (all amplitude on one of the four states)\\
    Hadamard $A$\\
    Hadamard $B$
\end{algo}

\item Explain why, given $n$ qubits, the operation that applies Hadamard to each qubit
puts equal amplitude on each of the $2^n$ basic states. You may provide either a formal proof or a less formal explanation.

\item Optional: watch the 3Blue1Brown video on quantum computing and Grover's algorithm: \url{https://www.youtube.com/watch?v=RQWpF2Gb-gU}.

\end{enumerate}



\end{document}
